% Scenario assumption, not an observed or measured before/after comparison.
\begingroup\setstretch{1}
\begin{tikzpicture}[
  font=\small,
  procbox/.style={draw=black!65, rounded corners=2pt, text width=2.9cm,
    minimum height=1.8cm, align=center, inner sep=4pt},
  flow/.style={-{Latex[length=2mm]}, thick},
  note/.style={align=left, anchor=west}
]
\node[note,font=\small\bfseries] at (0,1.3) {Angenommener Referenzprozess};
\node[procbox] (r1) at (1.6,0) {Auftrag auf\\ Papier erhalten};
\node[procbox,fill=black!12] (r2) at (5.4,0) {Ware entnehmen\\ und prüfen};
\node[procbox] (r3) at (9.2,0) {Menge, Mangel\\ auf Papier\\ festhalten};
\node[procbox,fill=blue!6] (r4) at (13,0) {Rückmeldung\\ später in Odoo\\ erfassen};
\draw[flow] (r1) -- (r2);
\draw[flow] (r2) -- (r3);
\draw[flow] (r3) -- (r4);
\node[note,font=\small\bfseries] at (0,-1.35) {Digital unterstützter Zielprozess};
\node[procbox,fill=blue!6] (s1) at (1.6,-2.7) {Auftrag in der\\ PWA öffnen};
\node[procbox,fill=black!12] (s2) at (5.4,-2.7) {Ware entnehmen\\ und prüfen};
\node[procbox,fill=blue!6] (s3) at (9.2,-2.7) {Pick bestätigen;\\ ggf. Mangel\\ digital erfassen};
\node[procbox,fill=blue!6] (s4) at (13,-2.7) {In Odoo\\ speichern und\\ Ergebnis zeigen};
\draw[flow] (s1) -- (s2);
\draw[flow] (s2) -- (s3);
\draw[flow] (s3) -- (s4);
\node[note,text width=14.5cm] at (0,-4.4)
  {Grau: Entnahme und Prüfung. Blau: digitale Erfassung oder Anzeige.\\
   Pfeile: vereinfachte Reihenfolge je Position; die Mängelmeldung ist optional.};
\end{tikzpicture}
\endgroup
